\documentclass[11pt,reqno]{amsart}

\usepackage[T1]{fontenc}
\usepackage{lmodern}
\usepackage{microtype}
\usepackage{mathtools,amssymb,amsfonts}
\usepackage{booktabs}
\usepackage{enumitem}
\usepackage[margin=1.08in]{geometry}
\usepackage{xcolor}
\definecolor{linkblue}{HTML}{244E73}
\definecolor{citegreen}{HTML}{315B45}
\usepackage[
  colorlinks=true,
  linkcolor=linkblue,
  citecolor=citegreen,
  urlcolor=linkblue,
  pdfauthor={Working draft},
  pdftitle={A counterexample to Makeev's regular-simplex conjecture in dimension four}
]{hyperref}
\usepackage[nameinlink,noabbrev]{cleveref}

\allowdisplaybreaks
\setlist[enumerate]{leftmargin=2.2em,itemsep=0.25em,topsep=0.4em}

\newtheorem{theorem}{Theorem}[section]
\newtheorem{proposition}[theorem]{Proposition}
\newtheorem{lemma}[theorem]{Lemma}
\theoremstyle{remark}
\newtheorem{remark}[theorem]{Remark}

\crefname{theorem}{theorem}{theorems}
\Crefname{theorem}{Theorem}{Theorems}
\crefname{lemma}{lemma}{lemmas}
\Crefname{lemma}{Lemma}{Lemmas}
\crefname{proposition}{proposition}{propositions}
\Crefname{proposition}{Proposition}{Propositions}
\crefname{remark}{remark}{remarks}
\Crefname{remark}{Remark}{Remarks}

\newcommand{\R}{\mathbb R}
\newcommand{\C}{\mathbb C}
\newcommand{\one}{\mathbf 1}
\newcommand{\SO}{\mathrm{SO}}
\newcommand{\eps}{\varepsilon}
\newcommand{\ip}[2]{\left\langle #1,#2\right\rangle}
\newcommand{\norm}[1]{\left\lVert #1\right\rVert}
\newcommand{\abs}[1]{\left\lvert #1\right\rvert}

\title[Counterexample in dimension four]{A counterexample to\\Makeev's
conjecture in dimension four}
\author{ChatGPT}
\thanks{This manuscript is AI-generated.}
\date{August 2026}

\begin{document}

\begin{abstract}
Let $u_{ij}=(e_i-e_j)/\sqrt2$ be the normalized roots of type $A_4$.
We give an explicit smooth odd function on $S^3$ which does not agree with
any linear functional on any rotated copy of these roots.  After identifying
$\R^4\cong\C^2$, for every sufficiently small nonzero real $\eps$ one may take
\[
 f_\eps(z_1,z_2)
 =\operatorname{Re}(z_1^2\overline z_2)
 +\eps\abs{z_1}^2\operatorname{Im}(z_1^2\overline z_2).
\]
The proof uses a cyclic model of the regular $4$-simplex, the exact-label
decomposition of edge flows, a cubic identity forced by a circle symmetry,
and an explicit algebraic transversality calculation.
\end{abstract}

\maketitle

\section{Introduction}

Kuperberg showed that a constant-width formulation of Makeev's universal-cover
conjecture is equivalent to asking whether every continuous odd function on
$S^{d-1}$ agrees with a linear function on a rotated $A_d$ root system.  He
proved the assertion for $d=3$ and showed that the homological obstruction used
there vanishes in higher dimensions \cite{Kuperberg1999}.  The result below gives a negative answer in the first higher
dimension.

The roots must be normalized because the function is defined on the unit
sphere.  Thus the standard model is
\[
 V=\one^\perp\subset\R^5,
 \qquad
 u_{ij}=\frac{e_i-e_j}{\sqrt2}\in S(V),
 \qquad 0\le i<j\le4.
\]
The configuration has an orientation-reversing symmetry, induced for example
by a transposition of two coordinate vectors.  Consequently the formulations
using $\mathrm O(4)$ and $\SO(4)$ are equivalent.

\begin{theorem}[Counterexample]\label{thm:counterexample}
There exists $\eps_0>0$ such that, for every real $\eps$ with
$0<\abs\eps<\eps_0$, the smooth odd function
\begin{equation}\label{eq:counterexample-function}
 f_\eps(z_1,z_2)
 =\operatorname{Re}(z_1^2\overline z_2)
 +\eps\abs{z_1}^2\operatorname{Im}(z_1^2\overline z_2),
 \qquad (z_1,z_2)\in S^3,
\end{equation}
has no $A\in\SO(4)$ and $b\in\R^4$ satisfying
\[
 f_\eps(Au_{ij})=b\cdot Au_{ij}
 \qquad(0\le i<j\le4).
\]
\end{theorem}

\section{The cyclic simplex and exact edge labels}\label{sec:setup}

Identify $\R^4$ with $\C^2$, equipped with the real inner product
\[
 \ip{z}{w}_{\R}
 =\operatorname{Re}(z_1\overline w_1+z_2\overline w_2).
\]
Let $\zeta=e^{2\pi i/5}$ and put
\begin{equation}\label{eq:cyclic-simplex}
 q_j=\sqrt{\frac25}(\zeta^j,\zeta^{2j}),
 \qquad j=0,\dots,4.
\end{equation}

\begin{lemma}[Cyclic regular simplex]\label{lem:cyclic-simplex}
The vectors $q_0,\dots,q_4$ satisfy
\[
 \sum_{j=0}^4q_j=0,
 \qquad
 \ip{q_i}{q_j}_{\R}
 =\delta_{ij}-\frac15.
\]
Consequently
\begin{equation}\label{eq:cyclic-roots}
 u_{ij}=\frac{q_i-q_j}{\sqrt2}
\end{equation}
are unit vectors forming an $A_4$ root system.  This system is
orientation-preservingly isometric to the normalized standard system.
\end{lemma}

\begin{proof}
Both $\sum_j\zeta^j$ and $\sum_j\zeta^{2j}$ vanish.  Moreover,
\[
 \ip{q_i}{q_j}_{\R}
 =\frac25\operatorname{Re}\bigl(\zeta^{i-j}+\zeta^{2(i-j)}\bigr).
\]
For $i\ne j$ the real part is $-1/2$, whereas for $i=j$ it is $2$.
Thus the displayed Gram matrix holds and $\norm{q_i-q_j}^2=2$.
The map $e_j-\frac15\one\mapsto q_j$ is therefore an isometry from the
standard simplex.  Composing it, if necessary, with an odd permutation of the
vertices makes it orientation-preserving without changing the unoriented root
system.
\end{proof}

By the last assertion, it is enough to prove \cref{thm:counterexample} for
the cyclic roots \eqref{eq:cyclic-roots}; we use the same notation $u_{ij}$
for them from now on.

For an odd function $F\colon S^3\to\R$ and $g\in\SO(4)$, define the skew
matrix
\begin{equation}\label{eq:label-matrix}
 Y_F(g)_{ij}=F(gu_{ij}),
 \qquad Y_F(g)_{ji}=-Y_F(g)_{ij}.
\end{equation}
Let
\[
 \Pi=I-\frac15\one\one^T.
\]

\begin{lemma}[Exact edge labels]\label{lem:exact-labels}
For a real skew $5\times5$ matrix $Y$, the following are equivalent:
\begin{enumerate}
\item there exists $c\in\R^4$ such that $Y_{ij}=c\cdot u_{ij}$;
\item there exist real numbers $\rho_0,\dots,\rho_4$ such that
      $Y_{ij}=\rho_i-\rho_j$;
\item $\Pi Y\Pi=0$.
\end{enumerate}
\end{lemma}

\begin{proof}
If $Y_{ij}=c\cdot u_{ij}$, take
$\rho_i=(c\cdot q_i)/\sqrt2$.  Conversely, subtract the mean of the
$\rho_i$ and put
\[
 c=\sqrt2\sum_i\rho_iq_i.
\]
The simplex Gram matrix gives $c\cdot q_j=\sqrt2\rho_j$, hence
$c\cdot u_{ij}=\rho_i-\rho_j$.

If $Y_{ij}=\rho_i-\rho_j$, then
$Y=\rho\one^T-\one\rho^T$, so $\Pi Y\Pi=0$.  Conversely, write
$J_0=\frac15\one\one^T$.  Since $Y$ is skew, $J_0YJ_0=0$, and
$\Pi Y\Pi=0$ gives $Y=J_0Y+YJ_0$.  With
$\rho_i=\frac15\sum_jY_{ij}$, this identity reads
$Y_{ij}=\rho_i-\rho_j$.
\end{proof}

This is the complete-graph cut/cycle decomposition, or equivalently the
one-dimensional part of combinatorial Hodge theory for edge flows
\cite{JiangLimYaoYe2011}.  We call such a matrix \emph{exact}.  The equality
in \cref{thm:counterexample} holds at a rotation $g$ precisely when
$Y_F(g)$ is exact: if $F(gu_{ij})=b\cdot gu_{ij}$, take $c=g^{-1}b$, and
conversely take $b=gc$.

We shall also use the elementary consequence that every skew matrix $X$ with
$X\one=0$ annihilates exact labels:
\begin{equation}\label{eq:cycle-annihilates-exact}
 \sum_{i<j}X_{ij}(\rho_i-\rho_j)=0.
\end{equation}

\section{A symmetry functional for the cubic}\label{sec:symmetry}

Put
\begin{equation}\label{eq:WQRH}
 W(z)=z_1^2\overline z_2,
 \qquad Q=\operatorname{Re}W,
 \qquad R=\operatorname{Im}W,
 \qquad H=\abs{z_1}^2R.
\end{equation}
Consider the circle action
\begin{equation}\label{eq:circle-action}
 h_\theta(z_1,z_2)=(e^{i\theta}z_1,e^{2i\theta}z_2),
 \qquad
 J(z_1,z_2)=(iz_1,2iz_2).
\end{equation}
The functions $Q$, $R$, and $H$ are invariant under $h_\theta$.

For $g\in\SO(4)$ write $p_i=gq_i$ and define
\begin{equation}\label{eq:X-Lambda}
 X(g)_{ij}=\ip{p_i}{Jp_j}_{\R},
 \qquad
 \Lambda_g(Y)=\sum_{i<j}X(g)_{ij}Y_{ij}.
\end{equation}
The matrix $X(g)$ is skew and
\[
 \sum_jX(g)_{ij}
 =\ip{p_i}{J\sum_jp_j}_{\R}=0.
\]
Thus \eqref{eq:cycle-annihilates-exact} shows that $\Lambda_g$ annihilates
every exact label matrix.

The following cubic identity is the only remnant of the variational cubic
argument needed for the counterexample.  Let
\[
 P_3(x)=\sum_{i=0}^4\ip{q_i}{x}_{\R}^3.
\]
For a homogeneous cubic $C$, write $T_C$ for its symmetric trilinear
polarization.  For any orthonormal basis $e_1,\dots,e_4$, set
\[
 \ip{C}{D}=\sum_{\alpha,\beta,\gamma=1}^4
 T_C(e_\alpha,e_\beta,e_\gamma)
 T_D(e_\alpha,e_\beta,e_\gamma).
\]
This is independent of the chosen orthonormal basis.  For
$K\in\mathfrak{so}(4)$, let
$(K\cdot C)(x)=\left.\frac{d}{dt}\right|_{t=0}C(e^{-tK}x)$.

\begin{lemma}[Cubic edge identity]\label{lem:cubic-edge-identity}
For every homogeneous cubic $C$ and $K\in\mathfrak{so}(4)$,
\begin{equation}\label{eq:cubic-edge-identity}
 \ip{C}{K\cdot P_3}
 =2\sqrt2\sum_{i<j}\ip{q_i}{Kq_j}_{\R}\,C(u_{ij}).
\end{equation}
\end{lemma}

\begin{proof}
The simplex is a Parseval frame: $\sum_iq_iq_i^T=I$.  Hence
\[
 \ip{C}{K\cdot P_3}
 =3\sum_iT_C(q_i,q_i,Kq_i).
\]
On the other hand, because $C(u_{ij})=C(q_i-q_j)/(2\sqrt2)$, the right side
of \eqref{eq:cubic-edge-identity} equals
\[
 \sum_{i<j}\ip{q_i}{Kq_j}_{\R}C(q_i-q_j).
\]
Expanding the cubic, the pure terms vanish since
$\sum_j\ip{q_i}{Kq_j}_{\R}=0$.  The two mixed terms combine to
\[
 -3\sum_{i,j}\ip{q_i}{Kq_j}_{\R}T_C(q_i,q_i,q_j).
\]
The Parseval identity gives
$Kq_i=-\sum_j\ip{q_i}{Kq_j}_{\R}q_j$, so the last expression is
$3\sum_iT_C(q_i,q_i,Kq_i)$, as required.
\end{proof}

\begin{lemma}[Global cubic annihilation]\label{lem:global-annihilation}
For every $g\in\SO(4)$,
\begin{equation}\label{eq:global-annihilation}
 \Lambda_g(Y_Q(g))=0.
\end{equation}
\end{lemma}

\begin{proof}
Take $K=g^{-1}Jg$ and $C=Q\circ g$ in
Lemma~\ref{lem:cubic-edge-identity}.  Then
\[
 \ip{q_i}{Kq_j}_{\R}=X(g)_{ij}.
\]
The orthogonal action on cubic tensors is infinitesimally skew-adjoint, while
$K\cdot(Q\circ g)=(J\cdot Q)\circ g=0$ by the circle invariance of $Q$.
Therefore
\[
 \ip{Q\circ g}{K\cdot P_3}
 =-\ip{K\cdot(Q\circ g)}{P_3}=0,
\]
and \eqref{eq:cubic-edge-identity} gives \eqref{eq:global-annihilation}.
\end{proof}

The remaining ingredient is an explicit algebraic statement.  Its proof is
contained in \cref{app:transversality}.

\begin{proposition}[Algebraic transversality]\label{prop:algebraic-transversality}
If $Y_Q(g)$ is exact, then
\begin{equation}\label{eq:transversality}
 \Lambda_g(Y_H(g))\ne0.
\end{equation}
\end{proposition}

\section{Completion of the counterexample}\label{sec:completion}

We now prove \cref{thm:counterexample}.  If the conclusion were false, there
would be nonzero numbers $\eps_\nu\to0$ and rotations $g_\nu\in\SO(4)$ for
which the labels of $Q+\eps_\nu H$ are exact.  By
Lemma~\ref{lem:exact-labels},
\begin{equation}\label{eq:defect-equation}
 \Pi Y_Q(g_\nu)\Pi
 +\eps_\nu\Pi Y_H(g_\nu)\Pi=0.
\end{equation}
After passing to a subsequence, compactness gives $g_\nu\to g$.  The second
term in \eqref{eq:defect-equation} tends to zero, so
$\Pi Y_Q(g)\Pi=0$; hence $Y_Q(g)$ is exact.

Because $\Lambda_{g_\nu}$ annihilates exact labels and
$\Lambda_{g_\nu}(Y_Q(g_\nu))=0$ by
Lemma~\ref{lem:global-annihilation}, exactness of
$Y_Q(g_\nu)+\eps_\nu Y_H(g_\nu)$ gives
\[
 \Lambda_{g_\nu}(Y_H(g_\nu))=0.
\]
Passing to the limit yields $\Lambda_g(Y_H(g))=0$.  This contradicts
Proposition~\ref{prop:algebraic-transversality}, because $Y_Q(g)$ is exact.  Therefore
an $\eps_0>0$ exists, and every function in
\eqref{eq:counterexample-function} with $0<\abs\eps<\eps_0$ is a
counterexample.  This proves \cref{thm:counterexample}.

\appendix

\section{Proof of algebraic transversality}\label{app:transversality}

We prove Proposition~\ref{prop:algebraic-transversality}.  Fix $g\in\SO(4)$ for which the
$Q$-labels are exact, write
\[
p_i=gq_i=(a_i,b_i)\in\C^2,
\qquad
a=(a_0,\dots,a_4),\quad b=(b_0,\dots,b_4)\in\C^5,
\]
and use coordinatewise multiplication of vectors in $\C^5$.  The Hermitian
inner product is conjugate-linear in the first variable:
\[
\ip{x}{y}_{\C^5}=\sum_i\overline{x_i}y_i.
\]

\subsection{Frame identities}

The $5\times4$ real matrix whose rows are the $p_i$ has orthonormal columns,
all perpendicular to $\one$.  Therefore
\begin{align}
\sum_i a_i=\sum_i b_i&=0,\label{eq:frame-sums}\\
\sum_i a_i^2=\sum_i b_i^2
=\sum_i a_ib_i=\sum_i a_i\overline b_i&=0,\label{eq:frame-orthog}\\
\sum_i\abs{a_i}^2=\sum_i\abs{b_i}^2&=2,\label{eq:frame-norms}\\
\abs{a_i}^2+\abs{b_i}^2&=\frac45.\label{eq:row-norm}
\end{align}
Consequently
\[
\one,a,b,\overline a,\overline b
\]
are mutually orthogonal in $\C^5$, with squared norms $5,2,2,2,2$.
Every $x\in\C^5$ has the expansion
\begin{equation}\label{eq:basis-expansion}
x=\frac{\sum_i x_i}{5}\one
+\frac12\left(
\ip{a}{x}a+\ip{b}{x}b
+\ip{\overline a}{x}\overline a
+\ip{\overline b}{x}\overline b
\right).
\end{equation}
For $u,v\in\C^5$, write
\[
(u\wedge v)_{ij}=u_iv_j-v_iu_j.
\]

\subsection{Cycle equations for the cubic}

At the root $gu_{ij}=(p_i-p_j)/\sqrt2$,
\[
W(gu_{ij})
=\frac1{2\sqrt2}(a_i-a_j)^2(\overline b_i-\overline b_j).
\]
Ignoring the harmless factor $1/(2\sqrt2)$, let
\[
\widetilde Y_{ij}=(a_i-a_j)^2(\overline b_i-\overline b_j).
\]
Modulo exact matrices, expansion gives
\begin{equation}\label{eq:cubic-cycle-form}
\Pi\widetilde Y\Pi
\equiv -a^2\wedge\overline b+2a\wedge(a\overline b).
\end{equation}
Indeed, the omitted term is
$a_i^2\overline b_i-a_j^2\overline b_j$.

Define
\begin{align}
s&=\frac12\sum_i\abs{a_i}^2a_i,
&\lambda&=\frac12\sum_i\abs{a_i}^2\overline b_i,\label{eq:s-lambda}\\
m&=\sum_i a_i^2\overline b_i,
&n&=\sum_i a_i^3,
&u&=\frac12\sum_i a_i\overline b_i^2,
&\ell&=\frac12\sum_i a_i^2b_i.\label{eq:m-n-u-ell}
\end{align}
Using \eqref{eq:basis-expansion},
\begin{align}
a^2&=sa+\frac m2b+\frac n2\overline a+\ell\overline b,
\label{eq:a2-expansion}\\
a\overline b&=\lambda a+ub+\frac m2\overline a-s\overline b.
\label{eq:abbar-expansion}
\end{align}
Substitution into \eqref{eq:cubic-cycle-form} yields
\begin{equation}\label{eq:omega}
\omega:=\Pi\widetilde Y\Pi
=2u\,a\wedge b
+m\,a\wedge\overline a
-3s\,a\wedge\overline b
-\frac m2b\wedge\overline b
-\frac n2\overline a\wedge\overline b.
\end{equation}
Exactness of the real $Q$-labels is precisely $\operatorname{Re}\omega=0$.
The six wedges of the complex basis $a,b,\overline a,\overline b$ are linearly
independent.  Comparing \eqref{eq:omega} with its conjugate gives
\begin{equation}\label{eq:cubic-exactness-conditions}
s=0,
\qquad m\in\R,
\qquad 4u=\overline n.
\end{equation}

Replacing $g$ by $h_\theta g$ changes neither the values of $Q$ and $H$ nor
the matrix $X(g)$, while
$\lambda\mapsto e^{-2i\theta}\lambda$.  We may therefore assume
\begin{equation}\label{eq:t-gauge}
\lambda=t\ge0.
\end{equation}
In addition define
\begin{equation}\label{eq:v-k}
v=\sum_i a_ib_i^2,
\qquad
k=\sum_i b_i^3.
\end{equation}

\subsection{The coordinatewise multiplication table}

Expanding each product in the orthogonal basis
$\one,a,b,\overline a,\overline b$ by \eqref{eq:basis-expansion}, and using
\eqref{eq:row-norm} and \eqref{eq:cubic-exactness-conditions}, gives
\begin{align}
a\overline a&=\frac25\one+tb+t\overline b,
&b\overline b&=\frac25\one-tb-t\overline b,\notag\\
a^2&=\frac m2b+\frac n2\overline a+\ell\overline b,
&\overline a^2&=\frac m2\overline b+\frac{\overline n}{2}a+\overline\ell b,
\notag\\
a\overline b&=ta+\frac{\overline n}{4}b+\frac m2\overline a,
&\overline a b&=t\overline a+\frac n4\overline b+\frac m2a,
\notag\\
ab&=ta+\ell\overline a+\frac v2\overline b,
&\overline a\,\overline b&=t\overline a+\overline\ell a
 +\frac{\overline v}{2}b,
\notag\\
b^2&=\frac n4a-tb+\frac v2\overline a+\frac k2\overline b,
&\overline b^2&=\frac{\overline n}{4}\overline a-t\overline b
 +\frac{\overline v}{2}a+\frac{\overline k}{2}b.
\label{eq:multiplication-table}
\end{align}
For example, the coefficient of $b$ in $a\overline a$ is
$\frac12\sum_i\overline b_i\abs{a_i}^2=t$, and the other coefficients are
obtained in the same way.

Coordinatewise multiplication is associative.  Comparing coefficients in
\eqref{eq:multiplication-table} gives the following necessary equations:
\begin{align}
3mn&=4tv,\tag{E1}\label{eq:E1}\\
\overline n\,v+16\ell t+8mt&=0,\tag{E2}\label{eq:E2}\\
-2km+24\ell t-n^2&=0,\tag{E3}\label{eq:E3}\\
-kt+\ell m+4t^2&=0,\tag{E4}\label{eq:E4}\\
\overline n\,k+mn+8tv&=0,\tag{E5}\label{eq:E5}\\
20\abs\ell^2+5m^2+5\abs n^2-40t^2&=8,\tag{E6}\label{eq:E6}\\
16\abs\ell^2-4m^2-\abs n^2+4\abs v^2&=0,\tag{E7}\label{eq:E7}\\
20\abs\ell^2+60t^2+5\abs v^2&=8.\tag{E8}\label{eq:E8}
\end{align}
More precisely, (E1)--(E3) are the $a,b,\overline b$ coefficients of
$(a^2)b=a(ab)$; (E4)--(E5) are the $a,b$ coefficients of
$(ab)b=a(b^2)$; and (E6), (E7), (E8) are respectively the indicated
coefficients of
\[
(a^2)\overline a=a(a\overline a),\qquad
(ab)\overline a=a(b\overline a),\qquad
(ab)\overline b=a(b\overline b).
\]

\subsection{Classification of the possible cubic zeros}

We extract the consequences of (E1)--(E8) needed below.

\medskip
\noindent\textbf{Case I: $t=0$.}
Equation (E1) gives $mn=0$.  If $m=0$, then (E3) gives $n=0$, while
(E6)--(E7) become
\[
20\abs\ell^2=8,
\qquad
16\abs\ell^2+4\abs v^2=0,
\]
a contradiction.  Hence $m\ne0$ and $n=0$.  Equations (E4) and (E3) give
$\ell=k=0$, and then
\begin{equation}\label{eq:case-I}
m^2=\frac85,
\qquad
\abs v^2=m^2.
\end{equation}

\medskip
\noindent\textbf{Case II: $t>0$ and $n=0$.}
Equation (E1) gives $v=0$, and (E2) gives
\begin{equation}\label{eq:case-II-ell}
\ell=-\frac m2.
\end{equation}
Here $m\ne0$, since otherwise \eqref{eq:multiplication-table} would give
$a^2=0$.  Equations (E3) and (E4) give
$k=-6t$ and $m^2=20t^2$.  Finally (E8) yields
\begin{equation}\label{eq:case-II}
t^2=\frac1{20},
\qquad
m^2=1.
\end{equation}

\medskip
\noindent\textbf{Case III: $t>0$ and $n\ne0$.}
First $m\ne0$, since otherwise (E1), (E2), and (E3) successively give
$v=0$, $\ell=0$, and $n=0$.  Put
\begin{equation}\label{eq:q-h}
q=\frac mt,
\qquad
h=\frac{\abs n^2}{t^2}>0.
\end{equation}
Equations (E1) and (E2) give
\begin{equation}\label{eq:v-ell-case-III}
v=\frac{3mn}{4t},
\qquad
\ell=-\frac{m(32+3h)}{64}.
\end{equation}
In particular, $\ell$ is real.  Equation (E5) gives
\[
k=-7m\frac n{\overline n}.
\]
Equation (E4) shows that $k$ is real, so
\begin{equation}\label{eq:sigma}
\sigma:=\frac n{\overline n}\in\{1,-1\},
\qquad
k=-7m\sigma,
\qquad
n^2=\sigma\abs n^2.
\end{equation}
After division by appropriate powers of $t$, equations (E4) and (E3) become
\begin{align}
3hq^2+32q^2-448\sigma q-256&=0,\label{eq:case-III-a}\\
-9hq-8\sigma h+112\sigma q^2-96q&=0.\label{eq:case-III-b}
\end{align}
Eliminating $h$ gives
\begin{align}
(q-4)(q+2)(21q^2+42q+16)&=0 &&(\sigma=1),\label{eq:resultant-plus}\\
(q+4)(q-2)(21q^2-42q+16)&=0 &&(\sigma=-1).
\label{eq:resultant-minus}
\end{align}
The roots $q=-2\sigma$ give $h=-64$ and are impossible.  One root of each
quadratic gives
$h=(-74-14\sqrt{105})/3<0$.  The candidates $q=4\sigma$ give $h=32$;
then \eqref{eq:v-ell-case-III} and (E1) imply
\[
\frac{\abs\ell^2}{t^2}=64,
\qquad
\frac{\abs v^2}{t^2}=288,
\]
and the left side of (E7), divided by $t^2$, is $2080$, a contradiction.
The only remaining possibility is
\begin{equation}\label{eq:case-III-qh}
q=-\sigma\left(1-\frac{\sqrt{105}}{21}\right),
\qquad
h=\frac{-74+14\sqrt{105}}3.
\end{equation}
Equation (E6) then gives
\begin{equation}\label{eq:case-III-tm}
t^2=\frac{7(5\sqrt{105}-33)}{1440},
\qquad
m^2=\frac{49\sqrt{105}-477}{1080}.
\end{equation}
These three cases exhaust every actual cubic zero, because every such zero
satisfies the necessary equations (E1)--(E8).

\subsection{Evaluation of the quintic functional}

For the actual root $gu_{ij}=(p_i-p_j)/\sqrt2$,
\begin{equation}\label{eq:quintic-edge}
\abs{z_1}^2W(z)
=\frac1{4\sqrt2}
(a_i-a_j)^3(\overline a_i-\overline a_j)
(\overline b_i-\overline b_j).
\end{equation}
Define
\begin{equation}\label{eq:Sigma}
\Sigma=\sum_{i<j}X(g)_{ij}
(a_i-a_j)^3(\overline a_i-\overline a_j)
(\overline b_i-\overline b_j).
\end{equation}
Then
\begin{equation}\label{eq:Lambda-H-Sigma}
\Lambda_g(Y_H(g))=\frac1{4\sqrt2}\operatorname{Im}\Sigma.
\end{equation}
In wedge notation,
\begin{equation}\label{eq:X-wedge}
X(g)=\frac i2\bigl(\overline a\wedge a+2\overline b\wedge b\bigr).
\end{equation}
Modulo an exact matrix, the degree-five edge expression in
\eqref{eq:Sigma} is
\begin{align}
F\equiv{}&
\overline b\wedge(a^3\overline a)
+\overline a\wedge(a^3\overline b)
-(\overline a\,\overline b)\wedge a^3\notag\\
&+3a\wedge(a^2\overline a\,\overline b)
-3(a\overline b)\wedge(a^2\overline a)\notag\\
&-3(a\overline a)\wedge(a^2\overline b)
+3(a\overline a\,\overline b)\wedge a^2.
\label{eq:F-wedge}
\end{align}
The omitted exact term is
$(a^3\overline a\,\overline b)\wedge\one$, which is annihilated by $X(g)$.
For the bilinear dot product $x\cdot y=\sum_i x_iy_i$,
\begin{equation}\label{eq:wedge-pairing}
\sum_{i<j}(u\wedge v)_{ij}(x\wedge y)_{ij}
=(u\cdot x)(v\cdot y)-(u\cdot y)(v\cdot x).
\end{equation}
Every dot product in \eqref{eq:wedge-pairing} is five times the coefficient of
$\one$ in the corresponding coordinatewise product.  Substituting
\eqref{eq:multiplication-table}, \eqref{eq:X-wedge}, and
\eqref{eq:F-wedge} gives
\begin{equation}\label{eq:Sigma-Braw}
\Sigma=-\frac i{80}\,\mathcal B_{\mathrm{raw}},
\end{equation}
where
\begin{align}
\mathcal B_{\mathrm{raw}}={}&
-960\overline k\,\ell t
+160\abs\ell^2m
+320\ell t^2
-120m^3\notag\\
&-75m\abs n^2
-800mt^2
+100m\abs v^2
+192m
-160nt\overline v.
\label{eq:Braw}
\end{align}
In all three cases above, $\ell$ is real.  The conjugates of (E4) and (E1)
give
\begin{equation}\label{eq:two-substitutions}
\overline k\,t=\ell m+4t^2,
\qquad
nt\overline v=\frac34m\abs n^2.
\end{equation}
Using \eqref{eq:two-substitutions}, followed by (E6) and (E8), reduces
\eqref{eq:Braw} to
\begin{equation}\label{eq:B-simplified}
\mathcal B
=m\bigl(180m^2+105\abs n^2-4400t^2-128\bigr)
-3520\ell t^2.
\end{equation}
The three cases now give
\begin{equation}\label{eq:B-values}
\begin{array}{c@{\qquad}c}
\toprule
\text{case}&\mathcal B\\
\midrule
\mathrm{I}&160m\\
\mathrm{II}&-80m\\
\mathrm{III}&60m(75-7\sqrt{105})\\
\bottomrule
\end{array}
\end{equation}
Each value is nonzero: in every case $m\ne0$, and
$75-7\sqrt{105}>0$.  Since $\mathcal B$ is real,
\eqref{eq:Lambda-H-Sigma} and \eqref{eq:Sigma-Braw} yield
\[
\Lambda_g(Y_H(g))=-\frac{\mathcal B}{320\sqrt2}\ne0.
\]
This proves Proposition~\ref{prop:algebraic-transversality}.


\begin{thebibliography}{9}

\bibitem{JiangLimYaoYe2011}
X.~Jiang, L.-H.~Lim, Y.~Yao, and Y.~Ye,
\emph{Statistical ranking and combinatorial Hodge theory},
Math. Program. \textbf{127} (2011), 203--244.
\href{https://doi.org/10.1007/s10107-010-0419-x}{doi:10.1007/s10107-010-0419-x}.

\bibitem{Kuperberg1999}
G.~Kuperberg,
\emph{Circumscribing constant-width bodies with polytopes},
New York J. Math. \textbf{5} (1999), 91--100.
\url{https://nyjm.albany.edu/j/1999/5-6p.pdf}.

\end{thebibliography}

\end{document}
